% SHARD-ID: LB-18-CONTINUUM-REDUCTION
% SHARD-TITLE: Reduction to the standard continuum setting
% SHARD-SOURCES: docs/reduction-limits.md (the working physics analysis);
%   theory/lanes/reduction/r1-smatrix.md, r3-wall.md, r2r4-ward.md
%   (independent codex calculation lanes, 2026-08-30; verdicts concur);
%   theory/lanes/reduction/q1-gauss.md, q2-memory-defs.md, q3-soft-defs.md,
%   q4-adversarial-defs.md (the DEFINITIONAL audit against 3+1 compact
%   Hamiltonian lattice QED, 2026-08-30; q4 hostile and run blind to q1-q3);
%   theory/verdicts/reduction-defs-adjudication-r1.md (adjudication);
%   docs/continuum-antecedents.md (digest, with one correction recorded here);
%   refs/arxiv-2201.01393 (type-B soft scaling, pole cancellation);
%   refs/arxiv-1602.08692 (magnon EFT, vanishing ferromagnetic scattering
%   length); refs/arxiv-1106.4382 (magnon-driven domain-wall motion);
%   refs/arxiv-1709.05018 (soft-pion triangle); refs/arxiv-1411.5745
%   (memory as a Fourier residue); refs/arxiv-1703.05448 (triangle logic);
%   definitions.md D6--D8, D12, D13--D18; claims/CLAIMS.md rows WI, A1, A2,
%   A2-orbit-r1, G0, G0-soft-r1, OR2, S2-2body, S2-2body-S, S-general, M,
%   M-flux, M-quant, M-quant-G, M-tk, Bc
% SHARD-CLAIMS: owns no statement.  Every result it uses is stated and proved
%   in the section that owns it; the reduction arguments themselves are
%   physical arguments and are deliberately outside claims/CLAIMS.md.

% --- shard-local semantic macros (rule 10 of the writing guide) -------------
\providecommand{\spindens}{M_{0}}          % spin (magnetisation) density S/a
\providecommand{\latconst}{a_{\ell}}       % lattice constant (a is used for windows)
\providecommand{\softexp}{\sigma}          % soft-scaling exponent of 2201.01393
\providecommand{\Gdiagr}{G_{\mathrm{diag}}}
\providecommand{\Aeffr}{\A_{\mathrm{eff}}}

\section{Reduction to the standard continuum setting}
\label{sec:continuum-reduction}

\subsection{The question}

Every physicist who reads this work will ask the same thing first, and they will
ask it before they ask anything else:

\begin{quote}
Does this reduce, in the right limits, to what we already accept as a memory
effect, a soft theorem, and an asymptotic symmetry?
\end{quote}

The question is fair and it is answerable. This section answers it in the
register in which it was asked: effective-field-theory matching, perturbation
theory, stationary phase, collective coordinates. Four reductions are examined
and each one ends in a verdict. Two of them come out well, one comes out badly,
and one comes out well in a way that requires a correction to be issued rather
than a confirmation to be announced. All four verdicts are stated plainly.

\begin{remark}[Status of everything in this section]
\label{rem:reduction-status}
The reduction arguments below are \emph{physical arguments, not part of the
claims DAG}. They carry no status tag, they are not proved anywhere in the
repository, and no result elsewhere in this labbook depends on them. Where a
reduction consumes a lattice result, that result is a theorem owned by another
section and is cited through its provenance block; where a reduction consumes a
continuum equation, that equation is transcribed from the local source. The
purpose of the section is to let a reader decide whether the lattice objects are
the objects they already know.
\end{remark}

\bigskip

\noindent Summary of the four verdicts.

\begin{center}
\begin{tabular}{@{}p{0.20\textwidth}p{0.74\textwidth}@{}}
\toprule
soft theorem & Reduces, with caveats. The naive continuum limit is a free
theory, but the soft \emph{law} survives it intact and matches the accepted
continuum magnon amplitude term by term. The soft-scaling exponent is
$\softexp=1$, not $2$. \\
\midrule
asymptotic symmetry & Does not reduce as stated. The naive one-dimensional
specialisation of the continuum construction is exactly the object this campaign
refuted. The charge-splitting half of the corner does reduce cleanly. \\
\midrule
memory & The magnonics reduction is clean, with constants. The Fourier-residue
reduction reduces in structure but not in content, for a reason the continuum
literature predicts. The kink-model transmission coefficient does not reduce at
all. \\
\midrule
the edges & Reduce, with caveats. Every silent continuum assumption has a named
lattice counterpart, and two of them the lattice removes outright rather than
renames. \\
\bottomrule
\end{tabular}
\end{center}

\subsection{The soft theorem}
\label{sec:reduction-soft}

\subsubsection*{What the lattice gives}

The lattice side of this reduction is exact. For the fully polarised bilinear
ferromagnet at any site spin $S$, the exact two-magnon coefficient ratio and its
soft phase slope are theorems (\cref{sec:two-magnon}): with $z_j=e^{ik_j}$,
$a=1+z_1z_2$, $b=z_1+z_2$ and $\mu=(2S-1)a+b$,
\begin{equation}
  S_{12}(k_1,k_2)=\frac{S\,a\,b-z_1\mu}{z_2\mu-S\,a\,b},
  \qquad
  \left.\frac{\partial\delta_{\mathrm{phys}}}{\partial k_s}\right|_{k_s=0}
  =\frac{\sgn(v_h-v_s)}{S}.
  \label{eq:red-exact-ratio}
\end{equation}
The soft magnon decouples in the plain sense that $S_{12}\to1$; this is a limit,
not an Adler-zero theorem.

\provenance{S2-2body-S, S2-2body, OR2, D6--D8}{\cref{sec:two-magnon};
\texttt{theory/spin-s-twomagnon.md}, \texttt{theory/oracle-bethe.md}}{---}{---}

\subsubsection*{The naive continuum limit is free, and this is verified}

Introduce a lattice constant $\latconst$ and physical momenta $p=k/\latconst$,
and let $\latconst\to0$ at fixed $p$. At $S=\tfrac12$ the exact phase is known
in closed form through the rapidity $\lambda(k)=\tfrac12\cot(k/2)$, and
\begin{equation}
  \lambda(\latconst p)=\frac{1}{\latconst p}-\frac{\latconst p}{12}+O(\latconst^3),
  \qquad
  \lambda_1-\lambda_2=\frac{p_2-p_1}{\latconst\,p_1p_2}+O(\latconst),
  \label{eq:red-rapidity}
\end{equation}
so that
\begin{equation}
  \delta_{12}
  =2\arctan\frac{1}{\lambda_1-\lambda_2}
  =2\arctan\frac{\latconst\,p_1p_2}{p_2-p_1}
  \;\xrightarrow[\ \latconst\to0\ ]{}\;0 .
  \label{eq:red-free}
\end{equation}
The rapidity difference diverges like $1/\latconst$ and the scattering matrix
tends to the identity. The naive continuum theory of lattice magnons is free.
This is the one-dimensional face of Dyson's vanishing ferromagnetic magnon
scattering length, and the continuum effective-theory literature reaches it
independently: in the ferromagnet without external fields the scattering length
vanishes, as a consequence of the derivative couplings demanded by Goldstone's
theorem.

\subsubsection*{What survives: the doubly soft amplitude}

The free limit is not the end of the story, because the soft theorem is not a
statement about the size of the interaction. Expanding
\eqref{eq:red-exact-ratio} to second order in \emph{both} lattice momenta, the
constant and linear parts of numerator and denominator cancel identically and
what remains is
\begin{equation}
  \boxed{\;
  S_{12}(k_1,k_2)=1-\frac{i}{S}\,\frac{k_1k_2}{k_1-k_2}+O(k^3).\;}
  \label{eq:red-doubly-soft}
\end{equation}
Two checks. Sending $k_1=k_s\to0$ at fixed $k_2=k_h>0$ returns the exact slope
$1/S$ of \eqref{eq:red-exact-ratio}. Setting $S=\tfrac12$ returns the frozen
spin-one-half coefficient $2$ of \cref{sec:two-magnon}.

In physical momenta, \eqref{eq:red-doubly-soft} reads
\begin{equation}
  S_{12}=1-\frac{i}{\spindens}\,\frac{p_1p_2}{p_1-p_2}+O(\latconst^2),
  \qquad
  \spindens:=\frac{S}{\latconst},
  \label{eq:red-doubly-soft-phys}
\end{equation}
where $\spindens$ is the spin, or magnetisation, density. Every microscopic
parameter has dropped out. The exchange coupling $J$ is absent, and the site
spin appears only inside $\spindens$.

\begin{observation}[The soft slope is an inverse spin density]
\label{obs:slope-is-inverse-density}
The exact lattice slope $\sgn(v_h-v_s)/S$ is a slope per \emph{lattice}
momentum. Per physical momentum it is
\begin{equation}
  \left.\frac{\partial\delta_{\mathrm{phys}}}{\partial p_s}\right|_{p_s=0}
  =\frac{\sgn(v_h-v_s)}{\spindens},
  \label{eq:red-slope-phys}
\end{equation}
and as a length --- the Wigner displacement that the soft magnon imparts to the
hard one --- it is $\latconst/S=1/\spindens$. The combination that stays finite
in the continuum is the spin density, and the soft law is the statement that the
displacement equals one inverse spin density.
\end{observation}

\noindent This resolves the apparent paradox of \eqref{eq:red-free}. A literal
$\latconst\to0$ at fixed $S$ sends $\spindens\to\infty$, that is, it describes an
infinitely stiff magnet, and of course the interaction vanishes there. In a real
magnet $\spindens$ is the saturation magnetisation, a finite measurable number,
and the magnon effective theory is a low-energy expansion with a physical
cutoff of order $1/\latconst$ rather than a continuum limit. The soft slope is
then a perfectly ordinary dimensionful Wilson coefficient.

\subsubsection*{Matching to the continuum magnon effective theory}

Model the magnon as a Schr\"odinger field with $\omega(p)=p^2/2m$. From
$\omega_S(k)=2JS(1-\cos k)\simeq JSk^2$ one reads $m=1/(2JS\latconst^2)$. In one
dimension the Born-level two-body scattering matrix for identical bosons is
$S-1=imM/\bigl(2\abs{p_1-p_2}\bigr)$ with $M$ the on-shell $2\to2$ amplitude, so
\eqref{eq:red-doubly-soft-phys} matches onto
\begin{equation}
  M(p_1,p_2)=-\frac{2p_1p_2}{m\,\spindens}=-4J\latconst^{3}\,p_1p_2 .
  \label{eq:red-matched-amplitude}
\end{equation}
Two consequences. First, writing the leading contact interactions as
\begin{equation}
  \mathcal{L}_{\mathrm{int}}
  =-\frac{g_0}{2}\bigl(\psi^\dagger\psi\bigr)^2
   -\frac{g_2}{4}\Bigl[\bigl(\partial_x\psi^\dagger\bigr)^2\psi^2+\text{h.c.}\Bigr]
   +\dotsb,
\end{equation}
the matching forces $g_0=0$ and $g_2\propto1/(m\spindens)=2J\latconst^3$: the
scattering-length operator has vanishing Wilson coefficient and the leading
magnon self-interaction carries two derivatives. That is Dyson's theorem,
obtained here as a consequence of \eqref{eq:red-doubly-soft} rather than as an
input. Second, the accepted continuum computation for the ferromagnet, with
dispersion $\omega=(F^2/\Sigma)\bm{k}^2$, gives the tree amplitude
\begin{equation}
  M\bigl[\pi(\bm{k}_1)+\pi(\bm{k}_2)\to\pi(\bm{k}_3)+\pi(\bm{k}_4)\bigr]
  =\frac{2F^2}{\Sigma^2}\,\bigl(\bm{k}_1\cdot\bm{k}_2\bigr),
  \label{eq:red-continuum-amplitude}
\end{equation}
stated in the source to agree with Dyson's microscopic analysis. The functional
form of \eqref{eq:red-matched-amplitude} and \eqref{eq:red-continuum-amplitude}
is the same, with no constant piece in either; identifying $\Sigma$ with the
spin density $\spindens$ and $F^2$ with the spin stiffness $JS^2\latconst$ makes
every parametric dependence agree, and leaves an overall factor of two that
depends on normalisation conventions this section has not re-derived. That
factor is recorded as an open point below.

\subsubsection*{Do the soft and continuum limits commute?}

They do, and it is worth saying why the question has a slightly disappointing
answer. The lattice constant enters the exact scattering matrix only through the
change of variables $k=\latconst p$, so there is nothing available to fail. On
the exact spin-one-half function \eqref{eq:red-free}, both orders of the two
limits give zero, and both orders of the two limits applied to the slope give
zero as well.

The limit that genuinely fails to commute is a different pair, and it is the one
that carries the physics.

\begin{observation}[Soft against hard, not soft against continuum]
\label{obs:soft-hard-noncommute}
From \eqref{eq:red-doubly-soft},
\begin{equation}
  \lim_{k_h\to0}\;
  \left.\frac{\partial\delta_{\mathrm{phys}}}{\partial k_s}\right|_{k_s=0}
  =\frac1S,
  \qquad
  \left.\frac{\partial}{\partial k_s}
  \Bigl[\lim_{k_h\to0}\delta_{\mathrm{phys}}\Bigr]\right|_{k_s=0}=0 .
  \label{eq:red-noncommute}
\end{equation}
The two orders differ by the entire content of the theorem.
\end{observation}

\noindent This is the same non-uniformity that \cref{sec:two-magnon} records
when it fixes the hard momentum away from the zone centre, and it is visible in
the quadratic coefficient $\abs{v_h}/\omega_h=\cot(k_h/2)$, which diverges as
$k_h\to0$. Its physical reading is that ``soft'' is a statement about the
\emph{ratio} $k_s/k_h=p_s/p_h$, which does not know about $\latconst$ at all.
That is exactly why the soft law survives a limit in which the scattering matrix
becomes trivial: the theorem constrains a scale-invariant ratio, and only the
overall size of the interaction carries the factor of $\latconst$.

\subsubsection*{Where the type-B bound sits, and a correction}

The accepted continuum route to a soft theorem for a non-relativistic Goldstone
runs through current conservation. With the broken current matrix element
written as
\begin{equation}
  \bra{\Omega}J^{\mu}(x)\ket{\theta(\bm p)}
  =e^{-ip\cdot x}\bigl[\,ip^{\mu}F_1(\abs{\bm p})+i\delta^{\mu0}F_2(\abs{\bm p})\,\bigr],
  \qquad
  \omega^2F_1+\omega F_2-\bm{p}^2F_1=0,
\end{equation}
the current matrix element between scattering states factorises on the Goldstone
pole,
\begin{equation}
  \bra{\beta}J^{\mu}(0)\ket{\alpha}
  =\frac{i}{p^0-\omega(\abs{\bm p})}
   \bra{\Omega}J^{\mu}(0)\ket{\theta(\bm p)}\braket{\beta+\theta(\bm p)|\alpha}
   +R^{\mu}(p),
\end{equation}
and current conservation cancels that pole exactly, leaving
\begin{equation}
  \braket{\beta+\theta(\bm p)|\alpha}
  =\left.\frac{p^0R_0(p)+p^rR_r(p)}{\omega(\abs{\bm p})F_1+F_2}\right|_{p^0=\omega(\abs{\bm p})} .
  \label{eq:red-nradler}
\end{equation}
The Adler zero then follows only under the additional assumption --- flagged in
the source as not following from standard polology --- that $R^\mu(p)$ is
non-singular as $\bm p\to0$ on shell. Counting powers in
\eqref{eq:red-nradler} with a soft-scaling exponent $\softexp$ defined by
$\mathcal{A}\propto\varepsilon^{\softexp}$ under $\bm p\mapsto\varepsilon\bm p$
gives the two bounds
\begin{equation}
  \softexp\ge\min(m,n+1)\ \ \text{(type $A_m$, $d\ge m+1$)},
  \qquad
  \softexp\ge\min(2m,n+1)\ \ \text{(type $B_{2m}$)},
  \label{eq:red-sigma-bounds}
\end{equation}
where $n$ is the degree of the theory's generalised \emph{spatial} shift
symmetry, that is, the degree of the polynomial
$\alpha(x)=\epsilon_{\nu_1\dots\nu_n}x^{\nu_1}\cdots x^{\nu_n}$ under which the
Goldstone field shifts.

The ferromagnet magnon is type $B_2$, so $m=1$. Its broken generators produce
only the ordinary constant shift of the Goldstone field, so $n=0$, and
\eqref{eq:red-sigma-bounds} gives $\softexp\ge\min(2,1)=1$. The exact lattice
amplitude \eqref{eq:red-doubly-soft} vanishes \emph{linearly} in either soft
momentum, so $\softexp=1$ exactly. The mechanism is transparent in
\eqref{eq:red-nradler}: with a quadratic dispersion the temporal piece
$p^0R_0=\omega R_0$ is already of second order, but the spatial piece
$p^rR_r$ is of first order unless a spatial shift symmetry of degree at least
one forces $R_r(0)=0$, and the Heisenberg ferromagnet has no such redundant
symmetry.

\begin{remark}[Correction to the campaign's reading of the type-B bound]
\label{rem:sigma-correction}
An earlier internal digest recorded the expectation that the lattice soft
theorem should realise the enhanced type-B value $\softexp=2$, and glossed the
type-B bound as needing no additional shift symmetry. Both readings are wrong.
Linear vanishing \emph{is} $\softexp=1$. The two type-$B_2$ theories that do
reach $\softexp=2$ in the source do so because they possess a linear spatial
generalised shift symmetry, that is $n=1$, whereupon
$\min(2m,n+1)=2$. What is genuinely special about type B is that $n=1$
\emph{suffices} to reach $\softexp=2$, whereas type $A_1$ with $n=1$ still
obtains only $\min(1,2)=1$. The lattice ferromagnet has $n=0$ and lands on
$\softexp=1$, in agreement with the bound and with the exact two-magnon
computation. No proved statement of this campaign changes; only the expectation
was mistaken.
\end{remark}

\begin{observation}[Verdict for the soft theorem]
\label{obs:verdict-soft}
\emph{Reduces, with caveats.} The exact lattice soft law reduces to the accepted
continuum statement --- a derivative-coupled type-$B_2$ Goldstone with vanishing
scattering length, tree amplitude proportional to $\bm k_1\cdot\bm k_2$, and an
ordinary Adler zero --- with the slope $1/S$ per lattice momentum becoming one
inverse spin density per physical momentum. The caveats are that the literal
continuum limit at fixed site spin is free and the physical statement therefore
requires a finite lattice constant; that the soft-scaling exponent is $1$ and
not $2$; and that an overall factor of two in the amplitude match is a
normalisation question left unchecked.
\end{observation}

\subsection{The asymptotic symmetry}
\label{sec:reduction-symmetry}

\subsubsection*{The specialisation the continuum construction makes}

A continuum asymptotic-symmetry construction assigns one charge per function on
the celestial sphere. In one spatial dimension the celestial sphere degenerates
to two points, left and right infinity, a function on it is a pair
$(g_L,g_R)$, and the diagonal subgroup acts trivially on the physical content of
the pair. The specialisation is therefore
\begin{equation}
  \A=(G_L\times G_R)/\Gdiagr,
  \qquad\text{complete invariant}\quad [\,g_Lg_R^{-1}\,].
  \label{eq:red-naive-A}
\end{equation}
This is the correct guess, and it is the group this campaign started from.

\subsubsection*{It is the object this campaign refuted}

The refutation is recorded in \cref{sec:corner-a}: the orbit of
\eqref{eq:red-naive-A} is \emph{not} the set of vacuum pairs, and
$[\,g_Lg_R^{-1}\,]$ is \emph{not} the complete invariant. What the lattice
supplies in its place are two different objects. In the unbroken case the orbit
is $\Aeffr=G/N_\alpha$, where $N_\alpha$ is the kernel of the virtual
representation $\rho_\alpha\colon G\to\mathrm{PGL}(\chi)$ carried by the bond
space; the stabiliser contains the diagonal and equals it exactly when
$N_\alpha$ is trivial. In the broken case, and under transitivity, vacuum pairs
form $(G/H_\alpha)\times(G/H_\alpha)$ with stabiliser $H_\alpha\times H_\beta$,
and the complete diagonal invariant is a \emph{double} coset in
$H_\alpha\backslash G/H_\alpha$.

\provenance{A2-orbit-r1 (refuted), A1(e), A2(e)}{\cref{sec:corner-a};
\texttt{theory/corner-a.md}, \texttt{theory/corner-a-kinks.md}}{\texttt{theory/checks/corner\_a\_check.py} C6--C8, C7, C11}{\texttt{theory/verdicts/corner-a-r3.md}}

\begin{observation}[Where the two pictures agree and where they part]
\label{obs:symmetry-agree-part}
When the virtual representation is faithful and the symmetry is completely
broken, $N_\alpha$ and $H_\alpha$ are trivial, both lattice classifiers collapse
to $G$, and \eqref{eq:red-naive-A} is recovered. The continuum answer is
therefore the lattice answer in the structureless case. The lattice adds three
things with no continuum counterpart: the finite kernel $N_\alpha$, produced by
a projective representation on a finite-dimensional bond space that the
continuum construction has no analogue of; the vacuum stabiliser $H_\alpha$,
which turns single cosets into double cosets; and the cohomology class
$[\omega_\alpha]\in H^2(G,U(1))$, an invariant of the \emph{representation}
rather than of the group, which obstructs removing the multiplier and has no
image in \eqref{eq:red-naive-A} at all. Conversely, the continuum has one piece
of machinery with no lattice image, the logarithmic celestial Goldstone
correlators, which rest on representation theory of the massless little group
and should not be sought here.
\end{observation}

\subsubsection*{The soft and hard charges do reduce}

The other half of the corner reduces cleanly, and it is the half a referee will
recognise immediately. The continuum construction defines a soft charge as the
zero-frequency antisymmetric combination of Goldstone creation and annihilation
operators, adds a hard charge built as a bilinear in the matter operators with
kernel $q^\mu p^\nu\gamma_{\mu\nu}/(p\cdot q)$, and proves that the sum
commutes with the scattering matrix between physical states. Its lattice mirror
is built from proved machinery.

The modulated charge $Q[f;\xi]=\sum_xf(x)q_x(\xi)$ obeys the exact continuity
identity
\begin{equation}
  \bigl[H,\,Q[f;\xi]\bigr]=\sum_x(\Delta f)(x)\,j_{x|x+1}(\xi)
  \label{eq:red-continuity}
\end{equation}
for every Lie-algebra element and every finite-range Hamiltonian, with the
second difference of the profile in the role of the derivative of the celestial
function; this is the lattice statement of current conservation and it involves
no boundary at infinity. Acting with $Q[f;\xi]$ on the vacuum creates a magnon
whose amplitude vanishes like the momentum, which is the soft charge; acting
with the same charge on a prepared hard magnon deforms the hard leg, and the
\emph{failure} of the soft charge to be absorbed is exactly the soft factor,
\begin{equation}
  Q_{k_s}\ket{k_h}-\ket{B^{\mathrm{in}}}
  =\bigl(1-S_{12}\bigr)\ket{P_{12}}
  =-2ik_s\ket{P_{12}}+O(k_s^2)
  \qquad\text{at }S=\tfrac12 .
  \label{eq:red-soft-hard-defect}
\end{equation}
This is the lattice form of ``soft charge plus hard charge commutes with the
scattering matrix'', with the exact slope of \cref{sec:two-magnon} in the role
of the continuum spin-rotation kernel. In both settings the soft factor is of
order one and not of order $1/\omega$: the lattice sits in the soft-pion
universality class, where the leading pole is absent, not in the photon or
graviton class.

\provenance{G0(e), G0(c), WI, AC-EX-2M-D29, S2-2body}{\cref{sec:symmetry-noether},
\cref{sec:corner-a}, \cref{sec:two-magnon}}{\texttt{theory/checks/corner\_a\_check.py} C3--C5, C9, C10}{---}

\begin{observation}[Verdict for the asymptotic symmetry]
\label{obs:verdict-symmetry}
\emph{Does not reduce as stated.} The naive one-dimensional specialisation of
the continuum asymptotic-symmetry group is literally the statement this campaign
refuted. The correct lattice classifiers reduce to it in the structureless case,
and the difference between them is precisely the structure the continuum
construction cannot see. The charge-splitting half of the corner --- continuity
equation, Ward identity, soft and hard charges --- reduces cleanly. For an
external presentation this means Corner A should be offered as a correction of
the naive specialisation, not as its confirmation.
\end{observation}

\subsection{The memory effect}
\label{sec:reduction-memory}

There are two different accepted continuum statements here, and the lattice must
be mapped onto both. They are not the same statement and they do not fare the
same way.

\subsubsection*{The magnonics reduction, with constants}

The accepted continuum result is that a spin wave passing a transverse domain
wall in an easy-axis ferromagnet drives the wall. The linearised
Landau--Lifshitz problem around the Walker profile becomes a Schr\"odinger
equation with a reflectionless P\"oschl--Teller potential,
\begin{equation}
  q^2\varphi(\xi)=\Bigl[-\frac{d^2}{d\xi^2}-2\,\mathrm{sech}^2\xi\Bigr]\varphi(\xi),
  \qquad
  \xi=\frac{z}{\Delta_W},\qquad q^2=\frac{\omega}{K}-1,
  \label{eq:red-pt}
\end{equation}
whose scattering solution has unit transmission and a pure phase shift. The
magnon's spin is $-\hbar$ on one side of the wall and $+\hbar$ on the other, so
each transmitted magnon deposits $2\hbar$ of angular momentum on the wall, and
the wall must move against the spin-wave direction at
\begin{equation}
  V_{\mathrm{DW}}=-\frac{\rho^2}{2}\,V_g,
  \qquad V_g=\frac{\partial\omega}{\partial k},
  \label{eq:red-ywx}
\end{equation}
with $\rho$ the dimensionless spin-wave amplitude.

The lattice side is the windowed memory theorem of
\cref{sec:memory-quantization}. With the frozen windowed wall coordinate
$\wallX_W$ and asymptotic densities $\pm s$, charge conservation gives the
displacement and the ledger
\begin{equation}
  \delta x=-\frac{\langle N_T\rangle}{s},
  \qquad
  2s\,\delta x+\bigl(q_{\mathrm{out}}-q_{\mathrm{in}}\bigr)=0,
  \qquad q_{\mathrm{in}}=-1,\quad q_{\mathrm{out}}=+1 .
  \label{eq:red-ledger}
\end{equation}

\provenance{M-quant, M-quant-G, M-flux, D13(a), D14, D17, D18}{\cref{sec:memory-quantization};
\texttt{theory/memory-quantization.md}, \texttt{theory/memory-quantization-general.md}}{\texttt{theory/checks/mquant\_check.py}, \texttt{theory/checks/mquant\_general\_check.py}}{corpus-r2, \texttt{theory/verdicts/mquant-g-r2.md}}

These are the same equation. The lattice arithmetic
$q_{\mathrm{out}}-q_{\mathrm{in}}=2$ \emph{is} the continuum statement that the
magnon spin flips from $-\hbar$ to $+\hbar$. To see it, write $\spindens=s/a_{\ell}$
for the spin density --- the same quantity that appeared in
\cref{obs:slope-is-inverse-density} --- and $\delta X=\latconst\,\delta x$ for
the physical displacement. Then $2s\,\delta x=2\spindens\,\delta X$ and
\eqref{eq:red-ledger} becomes
\begin{equation}
  2\spindens\,\delta X+\bigl(q_{\mathrm{out}}-q_{\mathrm{in}}\bigr)=0
  \quad\Longrightarrow\quad
  \delta X=-\frac{1}{\spindens}\ \ \text{per transmitted magnon},
  \label{eq:red-memory-phys}
\end{equation}
and, at a transmitted-magnon flux $\Phi$,
\begin{equation}
  V_{\mathrm{DW}}=-\frac{\Phi}{\spindens} .
  \label{eq:red-vdw}
\end{equation}
It remains to evaluate the flux for the continuum spin wave. The magnetisation
is a unit vector with transverse amplitude $\rho$ in the asymptotic domains, so
$m_z\simeq1-\rho^2/2$; each magnon removes one unit of spin, so the magnon
number density is $n=\spindens\rho^2/2$. The wall is reflectionless, so
$\Phi=nV_g$, and \eqref{eq:red-vdw} gives
\begin{equation}
  V_{\mathrm{DW}}=-\frac{\rho^2}{2}\,V_g,
\end{equation}
which is \eqref{eq:red-ywx} exactly, constants included, independently of the
exchange coupling, the anisotropy, the wall profile and the magnon momentum.

\begin{remark}[The two displacements are one inverse spin density]
\label{rem:two-twos-continuum}
The physical displacement per transmitted magnon in
\eqref{eq:red-memory-phys} and the Wigner displacement of
\cref{obs:slope-is-inverse-density} are numerically the same length,
$1/\spindens$. This gives the campaign's ``two twos'' cross-reference a
continuum reading: both quantities are one inverse spin density. It is an
observation about units and a physical argument only; the corresponding
conjecture is unaffected and remains a conjecture.
\end{remark}

\subsubsection*{The Fourier-residue chain}

The canonical continuum argument writes the memory as the residue of a
zero-frequency pole. Fourier transforming the radiative field by stationary
phase and assuming that it approaches finite but different values at early and
late retarded times, one obtains the memory as the coefficient of the pole in
frequency, and identifies that coefficient with the soft factor. In slogan form,
the Fourier transform of a pole in frequency space is a step function in time.

The residue \emph{structure} reduces correctly. The exact finite-time flux
identity of \cref{sec:memory-quantization} reads
\begin{equation}
  \delta x=\frac{1}{2s}\Bigl[\tilde{\jmath}_{a-1|a}(0)-\tilde{\jmath}_{b|b+1}(0)\Bigr],
  \label{eq:red-flux}
\end{equation}
which is precisely the zero-frequency weight of the physical boundary current,
obtained by exact telescoping of the windowed coordinate; and the spectral dress
of the memory observable writes the same number as a genuine direct-current
residue, $\delta x=\frac{1}{2s}\sum_{x\in W}\bigl[m_x(+\infty)-m_x(-\infty)\bigr]$.

What does \emph{not} reduce is the identification of that residue with the soft
factor. The continuum chain works because the product of frequency and radiative
field has a finite nonzero limit; there is a genuine pole, supplied by
Weinberg's soft factor. The lattice soft factor has no pole: the scattering
ratio tends to one and the soft-scaling exponent is $1$. Multiplying a regular
function by a vanishing frequency gives zero. The campaign's literal memory
conjecture --- displacement equals the zero-frequency limit of the soft factor
--- is refuted for exactly this reason, and what survives is the boundary-current
weight \eqref{eq:red-flux} together with the conditional quantisation
\eqref{eq:red-ledger}, into which soft data enter only through the transmission
probability.

\provenance{M (refuted), M-flux, D13(b)}{\cref{sec:memory-quantization};
\texttt{theory/memory-quantization.md} \S1}{\texttt{theory/checks/mquant\_check.py}}{corpus-r2}

This is not a lattice pathology. It is what the continuum literature predicts
for a derivative-coupled Goldstone. In the soft-pion triangle the leading soft
theorem is absent and the memory consequently lives in the \emph{subleading}
falloff coefficient of the pion field, where it equals a linear functional of
the escaped hard charge. The lattice reproduces that pattern exactly.

\begin{observation}[Which lattice observable is which falloff coefficient]
\label{obs:falloff-dictionary}
The leading radiative coefficient of the continuum field corresponds to the
outgoing one-magnon wave amplitudes on the two legs, that is to the reflection
and transmission amplitudes; these carry the phase and no direct-current shift.
The subleading coefficient, where the pion memory lives, corresponds to the
\emph{integrated} asymptotic tail content $\sum_{x>b}\bigl(\varrho(S^z_x)+s\bigr)=2sN_T$.
The continuum hard charge corresponds to $q_{\mathrm{out}}-q_{\mathrm{in}}=2N_T$,
and the continuum statement that the subleading coefficient equals a linear
functional of the hard charge corresponds to $\delta x=-\langle N_T\rangle/s$.
Both sides say the same thing: memory is a linear functional of the integrated
hard charge that escaped, not the residue of a soft pole.
\end{observation}

\begin{observation}[What plays the role of the falloff assumption]
\label{obs:falloff-hypotheses}
The continuum needs the radiative field to approach finite but different values
at early and late retarded times, and it needs the frequency and radius limits
taken in a definite order. The lattice counterparts are named hypotheses rather
than footnotes: the summability class of \cref{sec:local-relaxation}, which
requires the deviation of the magnetisation profile from its two asymptotic
values to be summable on each side, is the lattice ``finite but different
values''; the wave-operator and local-decay hypothesis is the lattice ``the
radiation has left''; and the spectral dress requires the time derivative of
each site magnetisation to be integrable in time. The lattice also states its
own order-of-limits fence in the same place, and it is the exact counterpart of
the continuum footnote: a plane-wave magnon is not summable, so every soft
statement about memory must fix the wave packet first and take the momentum to
zero afterwards, and the zero-frequency limit must be taken after the
thermodynamic limit.
\end{observation}

\subsubsection*{The kink model in its continuum scaling limit}

The easy-axis kink model of \cref{sec:kink-model} carries an unconditional
projected transmission coefficient,
\begin{equation}
  t(k)=\Bigl[1+\frac{iJ^2}{4\,\omega(k)\,v(k)}\Bigr]^{-1},
  \qquad
  T(k)=16(\Delta-1)^2k^2+O(k^4)\ \ \text{as }k\to0,
  \label{eq:red-tk}
\end{equation}
with crossover momentum $k_*=1/\bigl(4(\Delta-1)\bigr)$. The natural continuum
comparison is the calculation \eqref{eq:red-pt}. The comparison has one clean
half and one bad half.

The clean half is the dispersion. Writing the kink parameter
$q=\Delta-\sqrt{\Delta^2-1}$ and $\Delta=1+\epsilon$, the wall width in lattice
units is $\xi_{\mathrm{DW}}=1/\ln(1/q)\simeq1/\sqrt{2\epsilon}$, matching the
continuum $\Delta_W=\sqrt{A/K}$ up to a factor $\sqrt2$ under
$A\sim J\latconst^2$ and $K\sim J(\Delta-1)$; and in the scaling limit
$\Delta\to1^+$, $k\to0$ with $q_c:=k\,\xi_{\mathrm{DW}}$ fixed,
\begin{equation}
  \omega(k)=J(\Delta-\cos k)\simeq\omega_{\mathrm{gap}}\bigl(1+q_c^2\bigr),
  \label{eq:red-scaling-dispersion}
\end{equation}
which is exactly the continuum relation $q^2=\omega/K-1$. The lattice magnon
sector reduces cleanly to the continuum Landau--Lifshitz magnon on a wide wall.

The bad half is the transmission. In the same scaling limit
$v=J\sin k\simeq Jq_c\sqrt{2\epsilon}$ and $\omega\simeq J\epsilon(1+q_c^2)$, so
\begin{equation}
  \frac{J^2}{4\omega v}\simeq\frac{1}{4\sqrt2\,\epsilon^{3/2}q_c(1+q_c^2)}
  \;\longrightarrow\;\infty,
  \qquad
  T\simeq\Bigl[4\sqrt2\,\epsilon^{3/2}q_c(1+q_c^2)\Bigr]^{2}
  \;\longrightarrow\;0 .
  \label{eq:red-tk-scaling}
\end{equation}
The projected lattice wall becomes totally reflecting exactly where the
continuum wall is exactly reflectionless.

\begin{remark}[Why the transmission comparison fails, and what it costs]
\label{rem:tk-mismatch}
The two computations are controlled in opposite regimes. The kink at
$\Delta\to1^+$ is a wide object: its exact zero-energy product family has
$q\to1$, and in the magnetisation product basis it is a broad superposition over
configurations carrying many domain walls. The sector reduction that makes
\eqref{eq:red-tk} exact keeps only the three-wall component, and the measured
leakage into five-wall configurations grows like the inverse square of the
anisotropy, from below one per cent at $\Delta=8$ to about ten per cent at
$\Delta=2$; extrapolated to $\Delta\to1^+$ it is of order one. Physically, the
continuum reflectionlessness comes from the wall's translational zero mode
sitting exactly at threshold, which is what makes the potential in
\eqref{eq:red-pt} transparent; the truncated model does have a side-coupled kink
level, but its detuning stays finite as the momentum vanishes, so the threshold
cancellation is absent. At large anisotropy the two do agree: the lattice
reflection probability falls off like the inverse square of the anisotropy and
the wall is transparent for all momenta above the small crossover scale, which
is the continuum reflectionless answer. The cost is a presentational one and it
should be paid openly: \eqref{eq:red-tk} must not be offered as the lattice
version of magnon-driven wall motion. The memory law \eqref{eq:red-ledger} is
that, it is a conservation law, and it is independent of the transmission
probability; the transmission probability only sets the size of the effect,
through the expected transmitted number, in the Ising regime.
\end{remark}

\begin{observation}[Verdict for memory]
\label{obs:verdict-memory}
The magnonics reduction \emph{reduces cleanly}, with constants, and it is the
strongest reduction in this section; its one caveat is that the lattice
statement is conditional on an asymptotic-completeness hypothesis which the
continuum derivation replaces by a closed-form solution, so that the lattice
assumes what the continuum computes. The Fourier-residue reduction
\emph{reduces with caveats}: the direct-current structure and the falloff
hypotheses have exact counterparts, but the identification of the memory with
the residue of the soft factor does not reduce, and the campaign's refutation of
that identification is what the continuum literature predicts for a
derivative-coupled Goldstone. The transmission coefficient of the kink model
\emph{does not reduce}.
\end{observation}

\subsection{The edges}
\label{sec:reduction-edges}

\subsubsection*{From Ward identity to soft theorem}

The cleanest continuum precedent proves that the asymptotic-symmetry derivation
of a soft theorem and the textbook broken-Ward--Takahashi derivation are the
same derivation. Its recipe is short: take the Ward--Takahashi identity for the
broken current, Fourier transform to an on-shell massless momentum, apply the
reduction formula to the remaining fields, and take the soft limit; the contact
terms on the right-hand side drop out, one term on the left gives the amplitude
with one soft Goldstone, and the matter term gives the soft factor.

The lattice runs the same four steps, and the first two are theorems. The
continuity identity \eqref{eq:red-continuity} is the lattice conservation law
and needs no boundary at infinity, matching the local continuum programme rather
than a null-infinity construction. The truncated-symmetry identity of
\cref{sec:corner-a} telescopes the symmetry to exactly two boundary-bond
insertions, which is the lattice contact term; and the lattice version is finite
and computed, with both boundary terms retained, since discarding them is
demonstrably false at any finite window. The third step replaces the plane-wave
Fourier transform by momentum superpositions of window vectors and the reduction
formula by the wave-operator construction of \cref{sec:wave-operators}. The
fourth step is a theorem at two legs: on the separated-preparation subclass the
fixed-time readout has the exact first jet with phase slope $\sgn(v_h-v_s)/S$.

The caveat is the load-bearing one, and it is stated in
\cref{sec:triangle-status}: at $n$ legs this is the general soft conjecture, and
the hypothesis that supports it has no exhibited instance. The lattice therefore
reproduces the continuum pattern as a template, and as a theorem at two legs,
not as a theorem at $n$ legs.

\subsubsection*{From soft theorem to memory}

This edge does not reduce, and the reason has already been given: on the lattice
the soft factor has no pole, so the Fourier residue of the soft factor is zero,
and the corresponding conjecture is refuted. The surviving lattice edge runs
through charge conservation instead, which is the soft-pion pattern rather than
the graviton pattern. An external presentation may still use the slogan about
poles and step functions expositorily, provided it says in the same breath that
the lattice pole is absent and that the memory is therefore carried by
integrated charge rather than by a residue.

\subsubsection*{The silent assumptions, matched}

Every continuum derivation quoted in this section makes assumptions it does not
prove. The following correspondence records, as fact, which named lattice
hypothesis occupies each slot.

\begin{center}
\begin{tabular}{@{}p{0.44\textwidth}p{0.50\textwidth}@{}}
\toprule
\textbf{continuum assumption} & \textbf{lattice counterpart, with status} \\
\midrule
The remainder function is non-singular on shell as the momentum vanishes; the
source flags this as not following from standard polology. &
The remainder bounds of the uninstantiated reduction hypothesis, and the
continuous-extension clause of the universality class. The telescoping identity
is the mechanism intended to \emph{prove} rather than assume this. The
hypothesis has no exhibited instance; this is the campaign's advertised
technical advantage and it is not yet cashed. \\
\midrule
No singularities beyond the Goldstone pole; no bilinear current insertions on
external legs. &
The claim that the naive kinematic prefactor is a universal soft factor is
\statusRefuted: hard data can enter at first order through the current. The
lattice found the loophole rather than assuming it shut. \\
\midrule
The radiative field approaches finite but different values at early and late
retarded times. &
The summability class of \cref{sec:local-relaxation}, the wave-operator and
local-decay hypothesis, and integrability in time of each site magnetisation
rate. The local-decay hypothesis is open on the full chain. \\
\midrule
The product of frequency and radius is large when the limits are taken. &
The order-of-limits fence: fix the wave packet, then take the momentum to zero;
take the zero-frequency limit after the thermodynamic limit. Frozen in the
definitions. \\
\midrule
Stationary-phase mode expansion at large radius, with no zeroth-order term in
the falloff expansion. &
The wave-packet construction of \cref{sec:wave-operators}, with disjoint
velocity supports. Established as a construction; its interface with the
fixed-time readout has named gaps. \\
\midrule
No surface terms in the Ward identity. &
The truncated-symmetry identity keeps both boundary-bond insertions
\emph{explicitly}, and setting them to zero is recorded as false at finite
window. \statusProved, with the surface terms computed. \\
\midrule
A hyperbolic foliation with prescribed falloff near timelike infinity, needed to
make the hard charge a finite integral. &
The windowed wall coordinate is a bounded local observable unconditionally, so
the displacement exists on every state and at finite volume, with no asymptotic
hypothesis and no order-of-limits clause. This slot the lattice removes rather
than renames. \\
\midrule
A physical-mode condition and a transversality condition, imposed by hand. &
The admissible-profile conditions on the charge profile, with the type
discipline that forbids quoting a fixed-momentum identity under a decaying-profile
hypothesis. Frozen in the definitions. \\
\bottomrule
\end{tabular}
\end{center}

\begin{observation}[Verdict for the edges]
\label{obs:verdict-edges}
\emph{Reduce, with caveats.} Every silent continuum assumption has a named
lattice counterpart, and in two slots --- windowed memory and surface terms ---
the lattice eliminates the assumption rather than renaming it. The caveats are
that the symmetry-to-soft edge is a theorem only at two legs; that the
soft-to-memory edge does not run through a Fourier residue of the soft factor
and must not be advertised as doing so; and that the regularity slot, the one
this campaign claims as its chief contribution relative to the continuum, is
still occupied by a hypothesis with no exhibited instance.
\end{observation}

\subsection{The definitional audit: what the discrete and continuous
definitions say about each other}
\label{sec:reduction-defaudit}

The four reductions above ask whether the campaign's \emph{results} land on the
accepted continuum statements. They do not answer a sharper and prior question:
whether the campaign's \emph{definitions} denote the objects whose names they
borrow. That question was put separately, and it was put in the form that makes
it hardest to dodge --- compare each definition, clause by clause, with compact
Hamiltonian lattice electrodynamics in three spatial dimensions, which is a
regulator in which the Maxwell photon and the Weinberg pole demonstrably exist.
Four independent audits were run, one of them hostile and run without sight of
the other three. They agree on every overlapping item. This subsection records
what they found.

Nothing here changes a status. No lane asked for one and none was needed: the
statuses in the claims DAG were already scoped honestly. What the audit changes
is the vocabulary the campaign is entitled to use.

\begin{remark}[Register of this subsection]
\label{rem:defaudit-register}
As with the rest of this section, everything below is physical argument. Where
a lattice-gauge fact is standard it is used as such; where a continuum equation
appears it is transcribed from the local sources; where a limit is asserted it
is asserted, not proved. The audit's value is negative and definitional: it
tells the reader which sentences about this campaign would be false.
\end{remark}

\subsubsection*{The Gauss-law nucleus, and the exact statement it supports}

Fix a finite oriented cubic lattice of spacing $\latconst$. Each link carries
the rotor Hilbert space $L^2(U(1))$, with holonomy $U_\ell=e^{i\theta_\ell}$ and
integer electric flux $E_\ell=-i\,\partial/\partial\theta_\ell$, so that
\begin{equation}
  [E_\ell,U_{\ell'}]=\delta_{\ell\ell'}U_\ell,
  \qquad
  \spec E_\ell=\Z .
  \label{eq:da-rotor}
\end{equation}
Matter of integral charge sits on vertices with charge operator $\rho_x$. Write
$s_{x\ell}=+1$ if $x$ is the tail of $\ell$, $-1$ if it is the head, and $0$
otherwise. The Gauss operator and the physical subspace are
\begin{equation}
  G_x=(\operatorname{div}E)_x-\rho_x,
  \qquad
  (\operatorname{div}E)_x=\sum_\ell s_{x\ell}E_\ell,
  \qquad
  G_xP_{\mathrm{phys}}=0 .
  \label{eq:da-gauss}
\end{equation}

The first thing the audit establishes is which lattice-QED operator the
campaign's regional symmetry is allowed to be compared with, and it is not the
obvious one. The full local gauge transformation
$\Gamma[\varepsilon]=\exp\bigl(i\sum_x\varepsilon_xG_x\bigr)$ is a redundancy:
$\Gamma[\varepsilon]P_{\mathrm{phys}}=P_{\mathrm{phys}}$, so its orbit is a
point and it cannot be the nontrivial transplant of a truncated on-site
symmetry. The correct counterpart is the matter-only regional rotation
\begin{equation}
  W_R(\alpha)=\exp\Bigl(i\alpha\sum_{x\in R}\rho_x\Bigr),
  \label{eq:da-Wr}
\end{equation}
the gauged descendant of a global on-site matter symmetry. This distinction
carries the whole comparison.

With it, the match is exact and not merely suggestive. The finite-graph
divergence theorem is an operator identity, every interior link occurring once
with each sign,
\begin{equation}
  \sum_{x\in R}(\operatorname{div}E)_x=\sum_{\ell\in\partial R}\sigma_{R\ell}E_\ell,
  \qquad
  \sigma_{R\ell}=\pm1\ \text{as $\ell$ exits or enters $R$},
  \label{eq:da-divthm}
\end{equation}
and combining \eqref{eq:da-gauss} with \eqref{eq:da-divthm} on the physical
subspace gives
\begin{equation}
  \boxed{\;
  W_R(\alpha)P_{\mathrm{phys}}
  =\prod_{\ell\in\partial R}e^{\,i\alpha\sigma_{R\ell}E_\ell}\,P_{\mathrm{phys}} .\;}
  \label{eq:da-telescope}
\end{equation}
This is the gauge-theory avatar of the truncated-symmetry identity of
\cref{sec:corner-a}. In one spatial dimension the entering left link
contributes $e^{-i\alpha E}$ and the exiting right link contributes
$e^{+i\alpha E}$, which is exactly the orientation of the two virtual bond
insertions --- the inverse implementer on the left cut, the implementer on the
right cut. The dictionary is one line,
\begin{equation}
  V_\alpha\bigl(e^{i\alpha}\bigr)^{\sigma_{R\ell}}
  \;\longleftrightarrow\;
  e^{\,i\alpha\sigma_{R\ell}E_\ell},
  \label{eq:da-dictionary}
\end{equation}
and it is quantitative, not pictorial. Two differences must be stated with it:
the lattice-QED equality holds only after projection onto the Gauss-law
physical subspace, whereas the campaign identity is a tensor identity before any
local constraint is imposed; and the campaign identity starts from a genuine
global on-site symmetry, with no local constraint anywhere in sight.

For an angle-dependent profile the identity acquires the term that the
continuum calls soft. Summation by parts against a real profile
$\varepsilon_x$, with $(d\varepsilon)_\ell=\varepsilon_y-\varepsilon_x$ on an
internal link $\ell=(x\to y)$, gives the three operators
\begin{equation}
  Q^{\partial}_R[\varepsilon]=\sum_{\ell\in\partial R}\sigma_{R\ell}
      \varepsilon_{x_{\mathrm{in}}(\ell)}E_\ell,
  \qquad
  Q^{H}_R[\varepsilon]=\sum_{x\in R}\varepsilon_x\rho_x,
  \qquad
  Q^{\nabla}_R[\varepsilon]=\sum_{\ell\subset R}(d\varepsilon)_\ell E_\ell,
  \label{eq:da-three-charges}
\end{equation}
and Gauss's law then supplies the finite-spacing soft/hard dictionary
\begin{equation}
  \boxed{\;
  Q^{\partial}_R[\varepsilon]\,P_{\mathrm{phys}}
  =\bigl(Q^{H}_R[\varepsilon]+Q^{\nabla}_R[\varepsilon]\bigr)P_{\mathrm{phys}} .\;}
  \label{eq:da-soft-hard}
\end{equation}
For constant $\varepsilon$ the gradient term vanishes and
\eqref{eq:da-soft-hard} collapses to \eqref{eq:da-telescope}. The converse
statement is the honest one and it is worth displaying, because the naive
slogan is false: the modulated matter rotation on its own obeys
\begin{equation}
  e^{iQ^{H}_R[\varepsilon]}P_{\mathrm{phys}}
  =e^{\,i\bigl(Q^{\partial}_R[\varepsilon]-Q^{\nabla}_R[\varepsilon]\bigr)}P_{\mathrm{phys}},
  \label{eq:da-not-boundary-only}
\end{equation}
which is \emph{not} boundary-only when $d\varepsilon\neq0$. A claim that every
modulated on-site operation telescopes purely to the cut would be wrong.

\begin{observation}[The arrow to the accepted soft/hard split, and what it costs]
\label{obs:da-gauss-arrow}
Under a limit sequence that keeps the ultraviolet, thermodynamic and null
stages separate --- send the spacing to zero at fixed physical sphere radius
$R$, box size $L$ and retarded time $u$, then $L\to\infty$, then $R\to\infty$ at
fixed $u=t-R$, then $u\to-\infty$ --- the three operators of
\eqref{eq:da-three-charges} become
\begin{equation}
  Q^{\partial}\longrightarrow Q^{+}_{\varepsilon},
  \qquad
  Q^{H}\longrightarrow Q^{+}_{H},
  \qquad
  Q^{\nabla}\longrightarrow Q^{+}_{S},
  \label{eq:da-arrows}
\end{equation}
with $Q^+_\varepsilon=Q^+_S+Q^+_H$ the single sphere-flux charge of the
continuum construction, not a sphere flux plus a further hard part. Each arrow
is a physical argument resting on Riemann-sum control of the flux, the
prescribed falloffs, and the null-boundary assumptions; none is a theorem of
this campaign. Two fences travel with them. At finite spacing $Q^{\nabla}$
should be called a gradient term and not a soft-photon operator: softness needs
the null and zero-frequency limits that come later, and the zero mode is itself
defined as a frequency limit, $N_z=\lim_{\omega\to0}\int du\,F^{(0)}_{uz}
e^{i\omega u}$. And for massive charged matter the hard charge does not all
cross null infinity, so the decomposition is null hard plus timelike hard plus
soft, not the two null integrals of \eqref{eq:da-arrows}.
\end{observation}

\provenance{WI, A1, A2, D3, D4, D10}{\cref{sec:corner-a},
\cref{sec:symmetry-noether}; audit in
\texttt{theory/lanes/reduction/q1-gauss.md} \S1--\S3 and
\texttt{theory/lanes/reduction/q4-adversarial-defs.md} \S5}{---}{\texttt{theory/verdicts/reduction-defs-adjudication-r1.md} (N2)}

\subsubsection*{The two instantiations compared: same definitional scheme,
different theories}

A correction of record is required here (adjudication amendment r1a, owner
ruling 2026-08-30). The hostile lane originally graded the comparison below
as a \emph{failure of the definition} --- ``the bond-implementer package does
not transplant, and no limit sequence inside the definition repairs it.''
That grading was a category error, and the asymptotic-symmetry corner is in
fact the part of the campaign whose definitional correspondence with the
accepted continuum notion is \emph{immediate}. The accepted definition of an
asymptotic symmetry is a scheme, not an algebra: take the transformations
compatible with the asymptotic structure, quotient by those acting trivially
there, implement the survivors as boundary-supported operators, and read off
the charge algebra together with whatever central or projective extension it
carries. The campaign's symmetry corner \emph{is} this scheme, applied to a
spin chain --- and the Gauss-telescope nucleus above is the \emph{same}
scheme applied to the rotor theory, where it produces the boundary Gauss
algebra. The scheme commutes with the change of theory; the resulting
algebras differ because the theories differ. A spin chain does not contain a
photon's edge modes, and no definition of asymptotic symmetry should
manufacture one; in particular, asking a fixed finite-dimensional virtual
bond to carry the rotor relation $[E,U]=U$ is not a test any theory-side
definition could be expected to pass, because the limit in question is taken
in the \emph{theory} (put a rotor on each link), after which the scheme
delivers the correct algebra --- as the nucleus shows, exactly, at finite
spacing.

The table below therefore stands, reread: it catalogues how the two
\emph{theories'} asymptotic data differ under the one common scheme. These
are scope notes for anyone tempted to identify the spin-chain algebra with
the QED edge algebra --- a claim the campaign never made --- and not
failures of any definition.

\begin{center}
\begin{tabular}{@{}p{0.20\textwidth}p{0.74\textwidth}@{}}
\toprule
\textbf{what differs} & \textbf{how the two theories' asymptotic data differ} \\
\midrule
the quotient &
The campaign divides the symmetry group by
$N_\alpha=\{g:\rho_\alpha(g)\ \text{scalar}\}$, a state-level projective
kernel: a phase invisible on one bond space. Lattice QED divides
\emph{allowed} boundary transformations by those acting \emph{trivially on
physical data}. A constant phase can be scalar on one fixed charge sector and
still be physically measurable as a character across charged sectors, where it
enters the hard soft factor. The two quotients discard different things. \\
\midrule
locality &
The campaign residue is a virtual insertion defined only through padded windows
of a matrix-product state. The boundary flux
$\sum_{\ell\in\partial R}\sigma_{R\ell}E_\ell$ is a physical boundary-link
observable --- or, after a subregion factorisation choice, an edge-mode
generator. No canonical isomorphism identifies a $\chi$-dimensional virtual
space with $L^2(U(1))$ on a cut link. \\
\midrule
boundary multiplicity &
The campaign has one finite group at each of two ends of a one-dimensional cut.
A three-dimensional region has $O(R^2/\latconst^2)$ independent boundary
samples and tends to a function group on the sphere; the continuum has one
conserved charge for \emph{every} boundary function. That enlargement is not
contained in the definition as written, and the eventually-constant profile
class, which has only two ends, cannot encode an angular function at all. \\
\midrule
spectrum &
Compact electric flux is unbounded and integer valued, and $[E,U]=U$ shifts
every flux eigenvalue, which no finite spectrum admits --- so the two
theories' boundary algebras are simply non-isomorphic at any finite $\chi$.
This is theory content, not a definitional defect: the rotor algebra lives on
the gauge links of the rotor theory, and the same scheme produces it there
(the nucleus above), while on a spin chain it produces the finite virtual
algebra, which is the correct answer for that theory. \\
\midrule
the extra datum &
Abelian electric boundary charges commute; the finite electric boundary algebra
carries no multiplier. The campaign's load-bearing extra datum is precisely a
finite-dimensional projective multiplier, the cohomology class
$[\omega_\alpha]\in H^2(G,U(1))$. That class is not the Gauss-charge algebra of
this target, and the asymptotic-symmetry quotient does not manufacture it. \\
\bottomrule
\end{tabular}
\end{center}

\noindent Two smaller points complete the picture. The padding that makes the
window insertion map injective has no continuum falloff counterpart at all: at
fixed padding depth its physical thickness even vanishes as the spacing goes to
zero, and it implies neither the field falloffs nor smoothness of the boundary
profile. And the finite-total-variation profile condition is a
\emph{one}-dimensional total-variation condition; a radially constant,
angularly nonconstant profile in three dimensions generally fails its naive
analogue, the appropriate sufficient replacement being angular $C^2$ regularity
together with falloff-weighted bounds. What survives from the profile
definition is structural discipline --- a compactly supported profile gives a
bona fide operator, a non-decaying profile must be defined by a limiting action,
a distributional kernel is not itself an admissible operator --- and not
definitional equality.

\begin{observation}[Verdict for the symmetry corner (amended r1a)]
\label{obs:da-verdict-gauss}
The asymptotic-symmetry corner is the part of the campaign whose
correspondence with the accepted continuum notion is definitionally
immediate: the scheme --- boundary-supported implementers of the
transformations that survive the quotient by trivially-acting ones, with the
charge algebra and its extension read off --- is the accepted definition,
and it instantiates correctly on both theories. On the rotor theory it
yields, exactly at finite spacing, the Gauss nucleus: regional matter charge
equals oriented cut-link electric flux, with the same orientation as the
campaign's two bond insertions, and the weighted boundary charge splits into
a hard matter term and a gradient term that flow, under the ordered limits,
to the accepted hard and soft charges. This is a connection to accepted
continuum electrodynamics, not an analogy. The honest residual scope notes:
the two theories' asymptotic \emph{algebras} are different objects and must
not be identified (the table above); the one-dimensional profile classes need
their stated angular replacements in three dimensions; the modulated matter
charge supplies the hard Noether part only; and reaching the Weinberg pole
still requires a separate scattering and reduction argument.
\end{observation}

\subsubsection*{The memory verdict: correct rotor arithmetic, wrong observable}

The accepted electromagnetic datum is a real transverse one-form. In the
boundary gauge the local source gives the endpoint identity
\begin{equation}
  e^2\partial_zN=\int_{-\infty}^{\infty}du\,F^{(0)}_{uz}
  =A^{(0)}_z\big|_{\mathcal I^+_+}-A^{(0)}_z\big|_{\mathcal I^+_-},
  \label{eq:da-memory-field}
\end{equation}
and the operational readout attached to it is a probe kick or a relative phase:
integrating the Lorentz force through the radiation epoch gives
$m\,\Delta v_A=q_{\mathrm p}\int dt\,E_A=-q_{\mathrm p}\Delta A_A$. The lattice
discretisation of \eqref{eq:da-memory-field} is the \emph{link angle}: with
$\theta_\ell=\latconst e A_i+O(\latconst^2e)$ and a principal lift,
$\Delta\widetilde\theta_\ell/(\latconst e)\to\Delta A_i$ as $\latconst e\to0$,
the compact period $2\pi/(\latconst e)$ receding to infinity. That is the
object a discretisation of memory has to reach.

The campaign's circle-charge integrality clause reaches something else, and it
reaches it exactly. The compact rotor statement and the campaign clause are the
same theorem:
\begin{equation}
  e^{2\pi in_\kappa}=e^{2\pi i\kappa}I
  \iff
  \spec n_\kappa\subset\kappa+\Z
  \qquad\text{versus}\qquad
  e^{2\pi iS^z_x}=cI
  \iff
  \spec S^z_x\subset\kappa+\Z,\ c=e^{2\pi i\kappa},
  \label{eq:da-int-clause}
\end{equation}
clause for clause, including the affine offset and its role as a background or
twist label. A twisted rotor sector, a background electric flux, and the
campaign's projective on-site phase all produce the same coset structure, and
they cancel identically in a difference taken with the same normalisation. But
in lattice electrodynamics the statement belongs to the \emph{link electric
field}, conjugate to a gauge angle; in the campaign it belongs to an
\emph{on-site matter charge}. Compactness alone identifies neither with the
other, and the difference has a quantitative consequence: the link rotor has
unbounded spectrum, so $\latconst e\,n_\ell$ has vanishing level spacing and can
tend to a real canonical momentum, whereas a fixed finite on-site spectrum gives
$\latconst e\,S^z\to0$ and fills nothing. A real limit from on-site weights
would need the available weight range to grow --- a growing local
representation, or a collective many-site sum with occupied weights of order
$1/(\latconst e)$ --- and no such representation-scaling limit is stated.

What the protocol measures is therefore charge-transfer counting statistics. At
fixed window and reference the windowed charge has spectrum in one affine coset,
$\spec\Qhat_{W,c_0}\subset\kappa_{W,c_0}+\Z$ with
$\kappa_{W,c_0}\equiv L\kappa+s(a+b-1-2c_0)$ modulo the integers; time
evolution is an automorphism and preserves the spectrum; hence for two
sequential projective readouts
\begin{equation}
  \nu=q_--q_+\in\Z .
  \label{eq:da-tpm-integer}
\end{equation}
This is offset cancellation between two classical outcomes, not commutativity
of two operators and not a dynamical memory theorem. It is a legitimate
full-counting-statistics protocol; the lattice theorem that consumes it is
correct; the naming is what fails.

\begin{observation}[The free-Maxwell counterexample]
\label{obs:da-maxwell-counterexample}
Separate the asymptotically free radiative sector of a scattering solution from
the hard charged sector that fixes the constraint data. Free Maxwell already
carries infinitely many large-gauge symmetries, and one may choose data whose
radiation-zone field has a smooth finite-energy pulse with a nonzero endpoint
shift \eqref{eq:da-memory-field}. This is accepted electromagnetic soft memory,
measurable as a relative phase between adjacent charged probes. Now transplant
the protocol. Applied to the separated radiative sector, the windowed
\emph{matter} charge has no charge to measure: every spectral projection is the
zero-charge projection, every counting law is the point mass $\delta_0$, and the
protocol returns zero while the accepted memory is nonzero --- the transplant is
vacuous. Replacing the measured observable by the radiative memory coordinate
does not repair it: that is a field quadrature with continuous spectrum, so the
integrality clause fails outright and the support theorem disappears.
Substituting photon number to recover an integer spectrum changes both the
symmetry charge and the observable, and a sharp memory step is carried by an
infrared coherent cloud whose photon number need not be finite even when its
energy and its classical memory are, so the first-moment tightness clause can
fail too. The same obstruction is plainer still in gravity, where the memory
tensor varies continuously with the hard momenta and no integer on-site circle
spectrum can be its discretisation.
\end{observation}

\provenance{D13, D26, D27, M-flux, M-INDEX-fin, M-INDEX-spec}{\cref{sec:memory-quantization},
\cref{sec:memory-index}; audit in
\texttt{theory/lanes/reduction/q2-memory-defs.md} \S2--\S5 and
\texttt{theory/lanes/reduction/q4-adversarial-defs.md} \S2}{---}{\texttt{theory/verdicts/reduction-defs-adjudication-r1.md} (N1)}

Two further mismatches are worth displaying because each is a named, closable
gap rather than a vague worry. The first is the endpoint. A settled detector
record implies the protocol's time-averaged endpoint,
\begin{equation}
  \Bigl|\tfrac1T\!\int_T^{2T}\!f(t)\,dt-f_+\Bigr|
  \le\sup_{t\in[T,2T]}\abs{f(t)-f_+}\longrightarrow0,
  \label{eq:da-cesaro-forward}
\end{equation}
but the converse is false: for $f(t)=\cos\omega t$ the same average is bounded
by $2/(\abs{\omega}T)$ and tends to zero although $f$ never settles. The
protocol can therefore assign an endpoint to a persistent local oscillation for
which the accepted final reading does not exist. This is an infrared, long-time
issue; no refinement of the spatial discretisation touches it, and closing it
needs a Tauberian upgrade that is not proved. The second is conjugacy. If a
projective measurement resolves the integer flux sharply, then for the angle
unitary $U\ket{n}=\ket{n+1}$ and the associated dephasing map $\mathcal D$,
\begin{equation}
  \operatorname{tr}\bigl[\mathcal D(\rho)\,U\bigr]=0
  \qquad\text{for every trace-class }\rho:
  \label{eq:da-conjugate-kill}
\end{equation}
a sharp charge measurement destroys exactly the phase coherence one would need
to read an angle. The campaign's protocol is one step further removed still,
measuring on-site matter charge rather than link flux.

\begin{observation}[The missing bridge, and the name that survives]
\label{obs:da-memory-name}
To pass from a hard charge history to a field memory one must impose Maxwell's
constraint, specify early and late Coulombic data and boundary conditions, and
invert the angular differential operator. Different angular current profiles
can share one total escaped integer and yet have different
$\Delta A_A(z,\bar z)$; the map is not determined by the counting law. No such
reconstruction is present in any of the audited definitions or claims, and no
limit can manufacture it. This is the campaign's principal outstanding
definitional debt, and it is the object the background framing always said
memory should live in.

What does survive is a narrower and quite defensible name. The lattice
displacement law $\delta x=-\braket{N_T}/s$ reduces cleanly, constants
included, onto accepted magnon-driven domain-wall motion --- that reduction is
\cref{sec:reduction-memory} above, and it is the strongest reduction in this
section. So ``memory'' is justified in the magnonics sense, as a permanent
displacement of a real object produced by a passing wave. It is not justified
as generic electromagnetic or gravitational radiative memory, and the package
must not be advertised that way without the reconstruction theorem.
\end{observation}

\subsubsection*{The soft-insertion verdict: which definition is the accepted one}

Here the audit's finding is partly favourable and the favourable half should be
said first. The campaign's asymptotic-leg definition --- one additional
delta-normalised asymptotic excitation of unit leg weight, with the same
amputation applied to the amplitudes with and without it, and with an explicit
statement that this is \emph{not} a charge-created or current-created vector ---
is the accepted LSZ-type definition, transcribed into lattice-magnon language.
Its normalisation fence is what makes it so, and it is exact: with
$Q_k\ket{\Omega}=\sqrt{Z_\rho}\,\ket{k}$ one has
\begin{equation}
  \frac{\norm{Q_q\ket{h}}}{\sqrt N}=\sqrt{Z_\rho-2/N}\longrightarrow\sqrt{Z_\rho},
  \label{eq:da-leg-norm}
\end{equation}
so a charge-created leg becomes the asymptotic leg only after this conversion in
the reduction limit, and a common finite-state factor must never be used to set
the leg constant.

The fixed-time protocol datum is a different object. It is deliberately
pre-asymptotic: it applies the modulated charge to an already prepared hard
packet at finite time and reads a normalised projection ratio, declining to
define reduction exhaustiveness, wave operators, or on-shell matching. On the
spin-one-half ferromagnet the difference from the incoming scattering wave is
computed exactly,
\begin{equation}
  Q_{k_s}\ket{k_h}-\ket{B^{\mathrm{in}}}
  =\bigl(1-S_{12}\bigr)\ket{P_{12}}
  =-2ik_s\ket{P_{12}}+O(k_s^2),
  \label{eq:da-d6-defect}
\end{equation}
with relative size $\sqrt2\,\abs{k_s}\bigl(1+O(k_s)\bigr)$. Because the
protocol ties its soft momentum to its soft scale $\epsilon$, this is a
first-order discrepancy --- and the sought first jet is first order too, since
$S_{\mathrm{phys}}-1=2i\sgn(v_h-v_s)k_s+O(k_s^2)$. The two definitions can
therefore agree on the intercept and disagree on the jet, which is the entire
content of the theorem. This is why the unrestricted identification is left
open rather than assumed.

The order in which the two can be reconciled is fixed, and the protocol gets it
right. Writing the stages explicitly,
\begin{equation}
  \underbrace{N\to\infty}_{\text{thermodynamic}}
  \;\prec\;
  \underbrace{T\to\pm\infty}_{\text{scattering}}
  \;\prec\;
  \underbrace{\bigl(W\uparrow\Z,\ \sigma\downarrow0\bigr)}_{\text{window and packet width}}
  \;\prec\;
  \underbrace{\epsilon\downarrow0}_{\text{soft, last}} ,
  \label{eq:da-order}
\end{equation}
and the two definitions sit at two different points of that chain:
\begin{equation}
  \text{fixed-time datum}
  \;\xrightarrow[\ \text{separated-preparation class only}\ ]
              {\ \text{outer limit at fixed }\epsilon\ }\;
  \text{asymptotic-leg multiplier average}
  \;\xrightarrow{\ \epsilon\downarrow0\ }\;
  \text{soft theorem.}
  \label{eq:da-soft-diagram}
\end{equation}
On the separated-preparation class --- smooth compactly supported packets in one
regular primitive two-magnon channel, with disjoint velocity supports at each
fixed $\epsilon$, growing initial separation, shrinking hard width, a settling
time that grows, and the ring size chosen last --- a norm bridge of the form
\begin{equation}
  \bigl\|e^{-iH_ST}I_SF_R-I_SU_0(T)M_{S_{\mathrm{phys}}}F_R\bigr\|
  \le K_{M,\epsilon}\,s_M(f_\epsilon)\,s_M(g_\sigma)
     \Bigl[(1+R)^{3-M}+(1+d_\epsilon T-R)^{3-M}\Bigr]
  \label{eq:da-norm-bridge}
\end{equation}
makes every readout error vanish at each fixed $\epsilon$, and the protocol's
ratio equals the accepted on-shell multiplier average exactly,
\begin{equation}
  \mathcal A_*(\epsilon)=\int\bigl[S_{\mathrm{phys}}(k,h)-1\bigr]\,d\mu_{*,\epsilon}(k,h),
  \label{eq:da-zero-remainder}
\end{equation}
with \emph{zero} remainder after the outer limit --- so automatically $o(\epsilon)$
when the soft limit is finally taken. No bound uniform in $\epsilon$ is asserted
or needed. Outside that class no such theorem exists and
\eqref{eq:da-d6-defect} is the control. The gap between the two definitions is
therefore definitional, not merely technical.

\provenance{D24, D25, D29, AC-EX-2M-D29, S-IDX-MATCH-HS-SEP, ML6-HS-SEP}{\cref{sec:two-magnon},
\cref{sec:wave-operators}, \cref{sec:soft-index}, \cref{sec:universality};
audit in \texttt{theory/lanes/reduction/q3-soft-defs.md} \S2, \S5}{---}{\texttt{theory/verdicts/reduction-defs-adjudication-r1.md} (N4)}

\begin{observation}[The Weinberg-pole control experiment]
\label{obs:da-weinberg-control}
The audit ran the obvious hostile test: put a real gauge photon through the
campaign's soft-insertion definitions and see whether they notice. They do.
In lattice electrodynamics with a hard charged leg of dispersion $E(\bm p)$ and
velocity $\bm v(\bm p)$, the propagator adjacent to emission of a transverse
photon of momentum $\bm k$ and energy $\omega_\gamma$ has denominator
\begin{equation}
  \bigl(p^0+\omega_\gamma\bigr)^2-E(\bm p+\bm k)^2
  =2E(\bm p)\bigl[\omega_\gamma-\bm v(\bm p)\cdot\bm k\bigr]+O(\abs{k}^2),
  \label{eq:da-adjacent-prop}
\end{equation}
and the lattice Ward--Takahashi identity
$\widehat q_\mu\Gamma^\mu(p+q,p)=eQ\bigl[G^{-1}(p+q)-G^{-1}(p)\bigr]$ ties the
vertex numerator to the same derivative of the inverse propagator. Summing
attachments over outgoing and incoming legs gives
\begin{equation}
  \mathsf S_{\mathrm{gauge}}(q)
  =\sum_i\eta_i\,eQ_i\,\frac{p_i\cdot\varepsilon}{p_i\cdot q}
  =\frac{\mathsf S_{-1}(\hat k)}{\abs{k}}+O(1),
  \qquad \eta_i=\pm1\ \text{for outgoing/incoming},
  \label{eq:da-weinberg}
\end{equation}
the leading term being a \emph{Laurent} residue, not a Taylor derivative:
$\abs{k}\mathsf S_{\mathrm{gauge}}$ has a finite directional limit while
$\mathsf S_{\mathrm{gauge}}$ has no value or derivative at zero momentum. This
sharply violates the campaign's regularity hypotheses --- the twice-differentiable
soft map with zero intercept and finite contact first jet, and the conjectured
linear vanishing of the soft multiplier. The definitions therefore reject the
gauge amplitude, and they reject it for the right and observable reason, namely
its negative first Laurent power. The complementary half of the control is
equally sharp: a literal transplant of the campaign's modulated \emph{matter}
charge does not create the pole, and does not create the photon leg either,
because on the physical subspace it is the hard part of a large-gauge charge by
\eqref{eq:da-soft-hard} and the soft radiative part is missing. The correct
statement is thus narrow and true: the definitions distinguish the gauge and
global infrared classes by scope and by failed regularity, not by one unified
charge-insertion definition, and the absence of a $1/\omega$ pole in the
one-dimensional spin chain is consistent with, and does not exclude, the
standard pole in three spatial dimensions.
\end{observation}

\noindent One further distinction from the same lane deserves recording,
because the two vocabularies collide on the phrase ``pure gauge''. The
campaign's Goldstone analysis calls a direction pure gauge when its tensor is a
virtual coboundary $A_\alpha X-XA_\alpha$, which changes no local bulk quantity;
the continuum calls the large-gauge mode locally pure gauge while
\emph{retaining} its asymptotic endpoint data as a physical charge and vacuum
label. The exact statement common to both is only that a zero-momentum
coboundary changes no local bulk field or state but can leave boundary data.
The campaign's null direction is the analogue of a \emph{proper}, small gauge
redundancy; the soft photon is the boundary mode of a \emph{broken large} one.
Equating the two would reverse the broken and unbroken content of both
definitions. An exact analogue would need an enlarged asymptotic or edge phase
space in which the telescoping endpoint vectors are retained and charged rather
than removed, and the current definitions do not construct one.

\subsubsection*{The hostile lane's ranked findings}

The fourth audit was run adversarially and without sight of the other three. It
was asked to find rats, it found some, and it ranked them by severity. The
ranking is reproduced here in substance because the severities, not the prose,
are the useful part.

\begin{center}
\begin{tabular}{@{}p{0.12\textwidth}p{0.24\textwidth}p{0.56\textwidth}@{}}
\toprule
\textbf{severity} & \textbf{object} & \textbf{finding} \\
\midrule
fatal &
the no-contact universality class &
Source-dependent, and selected by conditions on the very amplitudes whose
factorisation is at issue, so that the membership clauses substantially restate
the desired conclusion. It has no proved microscopic member and it is not the
asymptotic-state domain the continuum theorems use. The conditional implication
built on it remains correct as an implication; it is not presently a nonvacuous
continuum-style soft theorem, and must never be described as one. \\
\midrule
withdrawn (r1a) &
the bond implementer as a QED boundary charge &
Originally graded fatal; reclassified as a category error by the r1a
amendment (owner ruling 2026-08-30). The finding tested the spin-chain
algebra against the job of being the QED edge algebra, which no one claimed;
the definitional \emph{scheme} of the symmetry corner is the accepted one
and instantiates correctly on both theories (the Gauss nucleus). What
survives of the finding is the scope note that the two theories' asymptotic
algebras are non-isomorphic and must not be identified. \\
\midrule
fatal &
the integrality-plus-counting package as generic memory &
It measures integer compact-charge counting statistics. A free radiative sector
can carry a nonzero soft endpoint shift while the matter-charge counting law is
$\delta_0$; substituting the field shift destroys integrality
(\cref{obs:da-maxwell-counterexample}). \\
\midrule
major &
the wall coordinate, charged-dressing masquerade &
The windowed observable tends to a soliton centre only for a fixed, centred core
profile. Its physical change is
$\Delta X+\Delta C_f+(\Delta\,\text{core charge})/(2\spindens)$, so a change of
charged core dressing at fixed geometric centre is indistinguishable from a
translation. The stated packet-outside and padding conditions do not exclude it;
a centring condition is owed. \\
\midrule
major &
the dynamical dress, undefined subtraction &
The displayed dynamical coordinate is exactly the conserved regularised total
magnetisation, so its displayed memory is identically zero --- the trap the
definition itself records. Invoking the asymptotic hypothesis cannot change the
value of a formula: a leg-subtracted replacement operator must first be defined,
and it is not. \\
\midrule
major &
the excitation ansatz called a ``particle'' &
The ansatz supplies a tangent-space quotient and a boundary-remainder limit, not
a spectral band, a resolvent pole, or a wave operator. Particle status is
legitimate only where exactness is supplied independently, by the exact-band
hypotheses or a separate spectral theorem. The audit checked and found no
currently proved row that crosses this gap silently. \\
\midrule
major &
integrality read as an asymptotic QED charge &
It captures the constant compact weight lattice. For a generic real angular
profile the weighted hard sum is not integer valued, and the radiative term is
not an on-site weight-lattice operator at all. \\
\midrule
minor &
``spectral dress'' as a name &
An exact Fourier rewrite of the finite-window response; the word can suggest a
radiative-field or soft-amplitude identification that the identity does not
supply. \\
\midrule
minor &
the projective phase in a gauge reading &
Ordinary compact rotor flux has trivial phase; a nontrivial offset needs
additional background or projective matter structure and should not be
presented as the generic gauge-field case. \\
\bottomrule
\end{tabular}
\end{center}

\begin{observation}[What survived the hostile hunt]
\label{obs:da-survivors}
The summary sentence, as amended at r1a, is the one worth remembering:
\emph{every definition is internally exact lattice mathematics; the
asymptotic-symmetry corner is definitionally the accepted scheme, correctly
instantiated; what fails are two continuum-facing identifications} --- the
counting package as generic radiative memory, and the no-contact class as a
continuum asymptotic-state domain. Specifically, the wall coordinate survives narrowly as a
centred rigid-soliton collective coordinate; the spectral dress survives as
exact endpoint-jump and Fourier-residue kinematics at fixed window with an
integrable rate; the integrality clause survives as the correct compact circle
arithmetic, matching a trivial-phase rotor flux; the counting protocol survives,
renamed, as a coherent conditional two-time-measurement full-counting-statistics
definition for lattice charge transfer; the locality, packet-norm, amputation
and zero-intercept clauses survive as useful analytic definitions, the
asymptotic-leg clause being the accepted reduction-type definition; the bond
implementer survives internally as a valid finite endpoint construction on
padded windows and states; and the truncated-symmetry identity survives as the
exact constant-profile Gauss telescope of \eqref{eq:da-telescope}. What does
not survive is the identification of the second group with the continuum objects
whose names they carry.
\end{observation}

\begin{observation}[Naming discipline]
\label{obs:da-naming}
The audit's operative conclusion is a vocabulary rule, and it binds every
external presentation of this campaign. The safe names are: a
\emph{collective-coordinate charge ledger} for the wall displacement and its
conservation law; a \emph{finite-window Fourier response} for the spectral
dress; \emph{compact-charge measurement statistics} for the integrality and
counting package; and \emph{conditional exact-band scattering} for the
asymptotic-state construction. ``Memory'' is defensible in the magnonics sense
only, as in \cref{obs:da-memory-name}. ``Asymptotic symmetry'' remains the
correct name for the symmetry corner without restriction (amendment r1a): the
corner is the accepted definitional scheme, correctly instantiated, and the
only rule there is not to identify the spin-chain algebra with a gauge
theory's edge algebra --- a claim never made. Absent the missing bridges, the
campaign may not describe the counting package as generic electromagnetic or
gravitational radiative memory, and may not present the conditional
factorisation implication as a nonvacuous continuum-style soft theorem. Three
bridges would lift the remaining restrictions, and only these three: a
reconstruction theorem tying the charge ledger to the field-side memory; an
exhibited microscopic member of the no-contact class; and a centring
amendment to the wall coordinate.
\end{observation}

\provenance{D4, D5, D12, D13, D24, D26, D27, ML5-B, D24-VAL, S-general}{\cref{sec:corner-a},
\cref{sec:kink-model}, \cref{sec:memory-index}, \cref{sec:universality};
audit in \texttt{theory/lanes/reduction/q4-adversarial-defs.md} \S6--\S7}{---}{\texttt{theory/verdicts/reduction-defs-adjudication-r1.md} \S1--\S4}

\subsection{What is still owed}

The reductions above rest in places on physical judgement rather than
computation, and those places are listed here rather than left implicit.

\begin{enumerate}
\item The amplitude match of \eqref{eq:red-matched-amplitude} against
      \eqref{eq:red-continuum-amplitude} agrees in form and in every parametric
      dependence and differs by an overall factor of two. That factor is
      believed to be a normalisation convention relating the amplitude to the
      scattering ratio for identical particles, together with the normalisation
      of the order-parameter density in the continuum source. It has not been
      checked.
\item The vanishing of the contact coupling is read off the leading doubly soft
      expansion. Higher orders in the lattice constant are not computed and could
      in principle generate a contact term at second order.
\item The assertion that the structureless case of the lattice classification is
      the correct continuum corner is a matching statement, not a theorem.
\item The magnon number density used to convert the continuum spin-wave
      amplitude into a magnon flux is the linear-spin-wave identification, exact
      to second order in the amplitude. The match of \eqref{eq:red-ywx} is
      therefore a leading-order match, as is the continuum derivation itself.
\item The lattice memory law is conditional on an asymptotic-completeness
      hypothesis that is open on the full chain. A referee may reasonably say
      that the lattice has assumed the physics and derived the bookkeeping.
\item The statement that the sector-reduction leakage is of order one at
      $\Delta\to1^+$ is an extrapolation from three measured anisotropies, not a
      bound.
\item The threshold explanation of the transmission mismatch in
      \cref{rem:tk-mismatch} is a physical reading. What is established is only
      that the two computations disagree in the limit.
\item The assignment of degree zero to the ferromagnet's generalised spatial
      shift symmetry in \cref{rem:sigma-correction} is read off the structure of
      the broken generators; no Lie-algebraic classification has been run to
      exclude a redundant symmetry.
\item The three missing bridges of \cref{obs:da-naming} are owed and none is
      started. A reconstruction theorem tying the charge ledger to the
      field-side memory is the largest of them, and until it exists the
      counting package may be called radiative memory in no sense at all.
\item Every arrow of \cref{obs:da-gauss-arrow} --- the Riemann-sum convergence
      of the boundary flux, the spacelike-to-null deformation, and the
      hard/soft split at null infinity --- is a matrix-element or classical
      statement. No operator-norm continuum limit of the unbounded flux is
      asserted, and the interacting scaling limit in which it would hold is not
      constructed.
\item The centring amendment owed to the wall coordinate is a change to a
      definition, not to a proof. Until it lands, every downstream use of that
      observable carries the charged-dressing ambiguity as an unstated
      hypothesis.
\end{enumerate}

\noindent Finally, the single most exposed point. The memory corner is proved
where the continuum computes, and computed where the continuum is proved: the
memory law, which reduces cleanly, is conditional on an open
asymptotic-completeness hypothesis, while the only unconditionally computed
quantity in the kink model is computed in a regime disjoint from the one in
which the continuum comparison lives. A reader who accepts the magnonics
reduction will ask next what the lattice computes that is neither assumed nor
restricted to a projection, and the honest answer today is the two-magnon soft
slope of \cref{sec:two-magnon}, and nothing else in the memory corner.
